\section{Posterior density levels and auxiliary projections}
\label{rlevels:section}\label{sec:r-density-levels}

A state can be informative because its posterior concentrates on a small
part of the sphere.  The two comparisons in this section measure that
concentration at dyadic radii.  In each comparison, a state is tested
after independent Gaussian observations have been supplied to the
analysis.  Replacing some of those rows by the current block makes the
state update a conditional Markov channel.  The task is to bound the
information cost of returning from the current rows to independent ones.

The two arguments use different refinements.  The first retains only
the level of the largest normalized ball mass and uses a net to estimate
the independent projection entropy.  The second retains both the first
maximizing radius and the density level; its covering estimate cancels
the radius directly.  These choices lead to different row counts and
to complementary exact-density estimates.

We use the measure, Gaussian law \(\gamma_t\), and entropy conventions
of Section~\ref{sec:r-experiments}.  Locally, write \(N=d-1\) for
the sphere dimension; this avoids confusing it with the actual-row
count \(n\) in the second comparison.  The row symbols have the
following roles, with their numerical values specified in each
subsection:
\begin{center}
\begin{tabular}{lcccc}
\toprule
Comparison & Total rows & Actual rows & Fresh rows & Moment \\
\midrule
One level & \(m\) & \(k\) & \(j=m-k\) & \(q\) \\
Two labels & \(k\) & \(n\) & \(l=k-n\) & \(q\) \\
\bottomrule
\end{tabular}
\end{center}
Constants \(C,C_0\) are positive and absolute and may increase.
All claims are for sufficiently large \(d\), with an absolute
threshold unless otherwise stated.

We use the local spherical-ball and projection-density results,
Lemmas~\ref{lem:r-caps} and \ref{lem:r-projection-density}.  In
particular, for every \(z\in\mathbb R^d\) and \(\rho>0\),
\begin{equation}\label{rlevels:ambient-cap}
 \sigma(B(z,\rho))\le (2\rho)^N.
\end{equation}
For \(\nu\le K\sigma\) and a full-row-rank matrix \(g\) with fewer
than \(N\) rows, write \(p_\nu(g,\cdot)\) for the jointly measurable
cube-average version of the density of \(gS\), \(S\sim\nu\).
For independent Gaussian \(G\) and \(S\sim\nu\), the logarithm
\(\log p_\nu(G,GS)\) is integrable.
The same lemma supplies the actual conditional image densities and
the integrable log ratios used below when a finite state is correlated
with the rows.  These density statements concern the exact
observations.

We use \eqref{eq:r-entropy-variational} with \(f=q\log F\) in the form
\begin{equation}\label{rlevels:tilt}
 \mathbb E_P\log F
 \le \frac{D(P\Vert Q)+\log\mathbb E_Q F^q}{q},
 \qquad P\ll Q,\quad q>0,\quad F\ge0.
\end{equation}
Truncation gives the integrable form and the upper bound when the
left side is negative infinite.

\subsection{A comparison using one density level}
\label{rlevels:one-level}

Set, locally in this subsection,
\begin{equation}\label{rlevels:one-parameters}
 m=\lfloor d/4\rfloor,\qquad k=\lfloor d/32\rfloor,\qquad
 j=m-k,\qquad q=\lfloor d/16\rfloor.
\end{equation}
The auxiliary projection has \(m\) rows, of which \(k\) can be actual
rows and \(j\) remain independent.  The integer \(q\) is the moment
order.  For large \(d\),
\[
 m<N,\qquad m-k-(q-1)\ge1.
\]

For a probability measure \(\nu\le K\sigma\), \(K<\infty\), put
\(r_i=2^{1-i}\), \(i\ge0\), and define
\begin{equation}\label{rlevels:one-maximal}
 L_\nu(s)=\max_{i\ge0} r_i^{-m}\nu(B(s,r_i)).
\end{equation}
The \(i=0\) term is \(2^{-m}\).  On the other hand,
\[
 r_i^{-m}\nu(B(s,r_i))\le K2^N r_i^{N-m},
\]
which tends to zero uniformly in \(s\).  Hence the maximum is
attained among at most \(C(1+d+\log K)\) initial scales, and
\[
 2^{-m}\le L_\nu(s)\le Ke^{Cd}.
\]
The ball-mass functions are Borel, by integration of the Borel
indicator \(\ind_{\{\|s-u\|\le r_i\}}\).  Thus \(L_\nu\) and the
first maximizing scale are Borel.  For \(a=2^b\), \(b\in\mathbb Z\),
let
\[
 F_a=\{s:a\le L_\nu(s)<2a\},\qquad w_a=\nu(F_a).
\]
Only \(C(1+d+\log K)\) levels can be nonempty, and a nonempty level
has \(a\ge2^{-m}\).

\begin{lemma}[Independent projection entropy from a net]
\label{rlevels:one-net}
If \(G\sim\gamma_m\) and \(S\sim\nu\le K\sigma\) are
independent, then
\begin{equation}\label{rlevels:one-net-bound}
 \mathbb E\log p_\nu(G,GS)
 \ge \mathbb E\log L_\nu(S)-Cd
       -\log\!\bigl(C(1+d+\log K)\bigr).
\end{equation}
\end{lemma}

\begin{proof}
At each relevant scale \(i\), take a maximal \(r_i\)-separated set
of centers on the sphere, and assign each point to a nearest center,
breaking ties in a fixed order.  Denote the cells and centers by
\(C_{ih}\) and \(c_{ih}\).  Each cell is contained in
\(B(c_{ih},r_i)\).  At most \(7^d\) centers lie within \(3r_i\)
of a fixed center: the ambient balls of radius \(r_i/2\) about
those centers have disjoint interiors and lie in a ball of radius
\(7r_i/2\).

Let \(b_{ih}\) be the sum of the \(\nu\)-masses of cells with centers
within \(3r_i\) of \(c_{ih}\).  For \(s\in C_{ih}\),
\(\nu(B(s,r_i))\le b_{ih}\), and
\(\sum_h b_{ih}\le7^d\).  Partition the sphere by its first
maximizing scale \(i(s)\) and its cell at that scale.  Write \(w_{ih}\)
for the masses of these classes, omitting zero-mass classes.  If
\(N_0\le C(1+d+\log K)\) is the number of relevant scales, Jensen
gives
\begin{equation}\label{rlevels:one-net-weight}
 \mathbb E\log\frac{\nu(B(S,r_{i(S)}))}
                       {w_{i(S),h(S)}}
 \le \log\sum_{i,h:w_{ih}>0}b_{ih}
 \le d\log7+\log N_0.
\end{equation}

In a class \((i,h)\), let \(\nu_{ih}\) be the conditional law of
\(S\).  Its points are within \(r_i\) of \(c_{ih}\).  Comparing
the density of \(GS\) with the Gaussian density of mean \(Gc_{ih}\)
and covariance \(r_i^2I_m\), nonnegative relative entropy yields
\[
 \mathbb E_{G,S\sim\nu_{ih}}\log p_{\nu_{ih}}(G,GS)
 \ge m\log(1/r_i)-Cm,
\]
because independence gives
\(\mathbb E\|G(S-c_{ih})\|^2\le mr_i^2\).
Also \(p_\nu\ge w_{ih}p_{\nu_{ih}}\) almost everywhere.
Average its logarithm under each class.  The identity
\(L_\nu(S)=r_{i(S)}^{-m}\nu(B(S,r_{i(S)}))\) then gives
\[
 \mathbb E\log p_\nu(G,GS)
 \ge \mathbb E\log L_\nu(S)
 -\mathbb E\log\frac{\nu(B(S,r_{i(S)}))}
                       {w_{i(S),h(S)}}-Cm.
\]
Apply \eqref{rlevels:one-net-weight}.  The refinement is finite,
and every class is dominated by a finite multiple of \(\sigma\);
Lemma~\ref{lem:r-projection-density} justifies all logarithmic
expectations.
\end{proof}

The next bound uses the same density level to control the projection
when the \(k\) rows will later be allowed to depend on the state.
The order of averaging is essential: the \(j\) independent rows are
averaged before taking a \(q\)-th moment in the other rows.

\begin{lemma}[A moment on one level]\label{rlevels:one-moment}
Let \(w_a>0\) and \(\mu=\nu(\,\cdot\,\mid F_a)\).  For independent
\(X\sim\gamma_k\), \(B\sim\gamma_j\), and \(G=(X;B)\),
\begin{equation}\label{rlevels:one-moment-bound}
 \left[\mathbb E_X
   \bigl(\mathbb E_B p_\mu(G,Gs)\bigr)^q\right]^{1/q}
 \le e^{Cd}\frac{a}{w_a}(1+d+\log K)
\end{equation}
for \(\mu\)-almost every \(s\).
\end{lemma}

\begin{proof}
For every ambient ball and \(t>0\),
\begin{equation}\label{rlevels:one-two-powers}
 \mu(B(z,t))\le C_0^d\frac{a}{w_a}t^m,
 \qquad
 \mu(B(z,t))\le C_0^d\frac{K}{w_a}t^N
\end{equation}
with one absolute \(C_0\).  For the first estimate, if the ball has
positive \(\mu\)-mass, choose \(s_0\in F_a\cap B(z,t)\).  For
\(t\le1\), a dyadic radius between \(2t\) and \(4t\) contains
\(B(z,t)\cap F_a\), and its \(\nu\)-mass is at most
\(2a(4t)^m\).  For \(t>1\), use total mass one and
\(a\ge2^{-m}\).  The second estimate follows from
\(\mu\le(K/w_a)\sigma\) and \eqref{rlevels:ambient-cap}.

Apply the shared/fresh moment theorem, Theorem~\ref{thm:r-mixed-moment},
with \(k\) shared rows, \(j\) fresh rows, total row count \(m\),
moment order \(q\), and
\[
 B_0=C_0^d a/w_a,\qquad D_0=C_0^d K/w_a.
\]
Its angular gap is \(m-k-(q-1)\ge1\), and its radial gap is
\(N-m>0\).  The inverse factor in its one-point integration is
exactly
\[
 \|u-s\|^{-j}\operatorname{dist}(u-s,V)^{-k},
 \qquad \dim V\le q-1.
\]
The ratio \(D_0/B_0=K/a\) does not contain \(w_a\), and
\(\log^+(K/a)\le\log K+m\log2\).  The theorem therefore gives the
right side of \eqref{rlevels:one-moment-bound} uniformly in the
cube-average scale, for every fixed \(s\).

For completeness, the passage to the exact density has the
almost-everywhere scope asserted here.  At \(Gs\) the cube averages
converge for almost every \((G,s)\) under
\(\gamma_m\otimes\mu\): condition on \(G\), use absolute
continuity of its image law and Lebesgue differentiation, and then
apply Fubini.  For \(\mu\)-almost every fixed \(s\), Fatou first
in \(B\) and then in \(X\) after taking the \(q\)-th power gives
\eqref{rlevels:one-moment-bound}.  This passage estimates the exact
label; it does not change the learner's observation.  The expansion
into independent points underlying the moment theorem is related to
the projection-moment method of Sharan, Sidford, and Valiant
\cite[Section~7, Lemma~12]{SSV2019}; the multirow estimate used here
is the local theorem with the displayed parameters.
\end{proof}

We now specify both experiments.  For a finite variable \(V\)
jointly distributed with \(S\sim\sigma\), define
\[
 \mathcal D(V)=I_{\mathrm{ind}}(S;V\mid G,GS),
 \qquad G\sim\gamma_m
 \text{ independent of }(S,V).
\]
In the actual experiment, \(X\sim\gamma_k\) is independent of
\(S\sim\sigma\), \(U\) is any finite variable jointly distributed
with \((S,X)\), and \(B\sim\gamma_j\) is independent of
\((S,X,U)\).  Put \(G=(X;B)\) and
\[
 \widetilde{\mathcal D}(U)=I_{\mathrm{act}}(S;U\mid G,GS).
\]
The same definition applies to a finite refinement \(V\) determined
by \((S,U)\).
Thus \(G\) and \(S\) are marginally independent in both experiments,
but \(G\) can depend on \(U\) in the actual one.  The independent
experiment uses the same \((S,U)\) marginal as the actual one.

\Needspace{6\baselineskip}
\begin{proposition}[One-level auxiliary-row comparison]
\label{rlevels:one-comparison}
For the preceding experiments and parameters,
\begin{equation}\label{rlevels:one-comparison-bound}
 \mathcal D(U)\le\widetilde{\mathcal D}(U)
 +C\left(d+\log(1+d+H(U))+\frac{H(U)}d\right).
\end{equation}
The proposition applies to every finite \(U\) with the specified
joint law.
\end{proposition}

\begin{proof}
First take \(V=U\), or a finite refinement of \(U\) determined by
\((S,U)\).  Let \(P\) be the actual law and \(Q\) the law with the
same \((S,V)\) marginal and an independent \(G\sim\gamma_m\).
Their \((G,S)\) marginals also agree.  Write
\(\nu_v=P_{S\mid V=v}=Q_{S\mid V=v}\).
Proposition~\ref{prop:r-information-comparison} gives the exact identity
\begin{equation}\label{rlevels:one-information-identity}
 \mathcal D(V)-\widetilde{\mathcal D}(V)
 =I_P(V;GS\mid G)-I_Q(V;GS\mid G)-I_P(S;G\mid V).
\end{equation}
Its density bound specializes to
\begin{equation}\label{rlevels:one-density-comparison}
 \mathcal D(V)-\widetilde{\mathcal D}(V)
 \le \mathbb E_P\log p_{\nu_V}(G,GS)
       -\mathbb E_Q\log p_{\nu_V}(G,GS).
\end{equation}
The negative term in \eqref{rlevels:one-information-identity} is
what pays for the changed conditional image density.  All quantities
are finite by Lemma~\ref{lem:r-projection-density} and
\(I_P(S;G\mid V)\le H(V)\).

For \(u\) of positive probability let \(K_u=P(U=u)^{-1}\), so
\(\nu_u\le K_u\sigma\).  Let \(J\) be the integer with
\(S\in F_{2^J}\) for the levels of \(L_{\nu_U}\), and put
\(V=(U,J)\).  The level count and Jensen give
\begin{equation}\label{rlevels:one-label-entropy}
 H(J\mid U)
 \le \mathbb E\log\!\bigl(C(1+d+\log K_U)\bigr)
 \le \log\!\bigl(C(1+d+H(U))\bigr).
\end{equation}
Under \(Q\), refinement increases the mean log projection density
by \(I_Q(J;GS\mid G,U)\ge0\).  Applying
Lemma~\ref{rlevels:one-net} to each \(\nu_u\) therefore yields
\begin{equation}\label{rlevels:one-independent-lower}
 \mathbb E_Q\log p_{\nu_V}(G,GS)
 \ge \mathbb E\log L_{\nu_U}(S)-Cd
       -\log\!\bigl(C(1+d+H(U))\bigr).
\end{equation}

For the actual upper bound, condition on \(S=s,V=v=(u,b)\) and
write \(a=2^b\).  Conditional on \((s,v,X)\), \(B\) is still
independent with law \(\gamma_j\).  Jensen in \(B\) bounds the
conditional mean of \(\log p_{\nu_v}(G,Gs)\) by the mean, under
\(P_{X\mid s,v}\), of
\(\log\mathbb E_Bp_{\nu_v}((X;B),(X;B)s)\).
Apply \eqref{rlevels:tilt} with reference \(\gamma_k\), order
\(q\), and Lemma~\ref{rlevels:one-moment}.  The conditional
divergences average to
\begin{equation}\label{rlevels:one-row-kl}
 \mathbb E D(P_{X\mid S,V}\Vert\gamma_k)
 =I_P(X;V\mid S)=I_P(X;U\mid S)\le H(U).
\end{equation}
The first equality uses marginal independence of \(X,S\); the
second uses that \(J\) is determined by \((S,U)\).  Thus the
refinement creates no additional conditional row information.
Finite average divergence also gives \(P_{X\mid S,V}\ll\gamma_k\)
almost everywhere.  The exact moment is needed only for
\(\nu_v\)-almost every \(s\), which is precisely the conditional
signal law at \(V=v\).

If \(w_a=\nu_u(F_a)\), the logarithm of the moment bound is
\[
 Cd+\log a+\log(1/w_a)+\log(1+d+\log K_u).
\]
Averaging, using
\(\mathbb E\log(1/w_a)=H(J\mid U)\), and applying Jensen once more,
we obtain
\begin{equation}\label{rlevels:one-actual-upper}
 \mathbb E_P\log p_{\nu_V}(G,GS)
 \le Cd+\mathbb E\log a+H(J\mid U)
       +\log(1+d+H(U))+\frac{H(U)}q.
\end{equation}
The log expectations and the truncated use of
\eqref{rlevels:tilt} are justified by the integrability already
established.

Finally, \(\log a\le\log L_{\nu_U}(S)\), so the local-density terms
in \eqref{rlevels:one-independent-lower} and
\eqref{rlevels:one-actual-upper} cancel.  Refinement gives
\[
 \mathcal D(U)\le\mathcal D(V),\qquad
 \widetilde{\mathcal D}(V)
 \le\widetilde{\mathcal D}(U)+H(J\mid U).
\]
Combine these inequalities with \eqref{rlevels:one-density-comparison}
and \eqref{rlevels:one-label-entropy}.  Since \(q\) is a fixed
positive fraction of \(d\) for large \(d\), the result is
\eqref{rlevels:one-comparison-bound}.
\end{proof}

\subsection{A comparison retaining the radius and the level}
\label{rlevels:two-level}

We now use more auxiliary rows and retain the first maximizing radius
as a second label.  This lets a covering argument estimate the
independent entropy without the net-weight average above.  In this
subsection the local parameters are
\begin{equation}\label{rlevels:two-parameters}
 k=\lfloor d/2\rfloor,\qquad n=\lfloor d/8\rfloor,\qquad
 l=k-n,\qquad q=n.
\end{equation}
Here \(n\) is the number of actual rows; the surface exponent remains
\(N=d-1\).  In particular
\[
 N-k\ge d/3,\qquad k-q-n\ge1,\qquad N-q-n\ge1
\]
for large \(d\).

For a finite variable \(V\) with \(S\sim\sigma\), let
\[
 \mathcal J(V)=I_{\mathrm{ind}}(S;V\mid G,GS),
 \qquad G\sim\gamma_k
 \text{ independent of }(S,V).
\]
In the actual experiment \(A\sim\gamma_n\) is independent of
\(S\sim\sigma\), \(V\) has any joint law with \((S,A)\), and
\(B\sim\gamma_l\) is independent of \((S,A,V)\).  Put
\(D=(A;B)\).  Thus \(D\) is marginally independent of \(S\)
with law \(\gamma_k\), although its conditional law given
\((S,V)\) need not be Gaussian.  The independent experiment defining
\(\mathcal J(V)\) preserves the actual \((S,V)\) marginal.

\begin{proposition}[Two-label auxiliary-row comparison]
\label{rlevels:two-comparison}
If \(V\) is finite and \(H(V)\le d^2\) in the preceding actual
experiment, then
\begin{equation}\label{rlevels:two-comparison-bound}
 \mathcal J(V)\le I_{\mathrm{act}}(S;V\mid D,DS)+Cd.
\end{equation}
\end{proposition}

\begin{proof}
Let \(P\) denote the actual law.  We first construct the refinement and obtain the two geometric
estimates it supplies.  For \(\pi_v=P(V=v)>0\), write
\(\nu_v=P_{S\mid V=v}\le\pi_v^{-1}\sigma\), and define
\[
 \lambda_v(s)=\max_{i\ge0} r_i^{-k}\nu_v(B(s,r_i)),
 \qquad r_i=2^{1-i}.
\]
Let \(j\) be the first maximizing index and let
\(b=\lfloor\log_2\lambda_v(s)\rfloor\).  The maximum is attained:
its \(i=0\) term is \(2^{-k}\), while
\(\pi_v^{-1}2^N r_i^{N-k}\) bounds the \(i\)-th term and tends
to zero.  The ball-mass functions are Borel, so both labels are
measurable.  The same bounds give
\[
 b\ge-k,\qquad 2^b\le\pi_v^{-1}C^d,\qquad
 2^{-k}\le\pi_v^{-1}C^d r_j^{N-k}.
\]
Because \(N-k\ge d/3\), the last inequality implies
\(j\le C(1+\log(1/\pi_v)/d)\).  At fixed \(v\), the number of
pairs \((j,b)\) is at most \(C(d+1+\log(1/\pi_v))^2\).
For \(W=(V,j,b)\) evaluated at \(s=S\), this proves finiteness and
\begin{equation}\label{rlevels:two-label-entropy}
 H(W\mid V)\le C+2\log(d+1+H(V)).
\end{equation}

For \(w=(v,j,b)\) with \(\pi_w=P(W=w)>0\), set \(R=r_j\).
Let \(E_w\) be the corresponding subset of the sphere, and write
\[
 \nu_w=P_{S\mid W=w}
       =\alpha_w\,\nu_v|_{E_w},\qquad
 \alpha_w=\frac{\pi_v}{\pi_w}.
\]
Every \(s_0\in E_w\) satisfies
\[
 \nu_v(B(s_0,R))\ge2^bR^k,\qquad
 \nu_v(B(s_0,r_i))\le2^{b+1}r_i^k\quad(i\ge0).
\]
Choose centers in \(E_w\) successively at distances greater than
\(2R\) from those already chosen, until their \(2R\)-balls cover
\(E_w\).  Their \(R\)-balls are disjoint and each has
\(\nu_v\)-mass at least \(2^bR^k\); hence the process stops with
at most \((2^bR^k)^{-1}\) centers.  Assign each point to one of
its covering balls.  Under the independent matrix \(G\), conditioning
also on this finite index adds at most
\(-b\log2-k\log R\) to the entropy.  Within a ball, comparison with
a Gaussian of covariance \((2R)^2I_k\) adds at most
\(k\log R+Ck\), since
\(\mathbb E\|G(S-z)\|^2\le4kR^2\) for its center \(z\).
Consequently
\begin{equation}\label{rlevels:two-independent-entropy}
 h_{\mathrm{ind}}(GS\mid G,W=w)\le-b\log2+Ck.
\end{equation}
The entropies are finite by Lemma~\ref{lem:r-projection-density}.

The same labels give, for all \(z\in\mathbb R^d\) and \(\rho>0\),
\begin{equation}\label{rlevels:two-two-powers}
 \nu_w(B(z,\rho))
 \le C_0^d\min\{U_w\rho^k,\ \Gamma_w\rho^N\},
 \qquad U_w=\alpha_w2^b,\quad \Gamma_w=\pi_w^{-1}.
\end{equation}
For the first bound, choose a point of \(E_w\) in the ball when
there is one, double the radius, and use a dyadic radius in
\([2\rho,4\rho]\) if \(\rho\le1\).  The upper ball bound at that
point and restriction by \(\alpha_w\) prove the assertion.  For
larger \(\rho\), use total mass one and \(U_w\ge2^{-k}\).
The second bound follows from \(\nu_w\le\Gamma_w\sigma\) and
\eqref{rlevels:ambient-cap}.  Notice that \(\alpha_w\ge1\),
\(\Gamma_w\ge1\), and \(U_w\ge2^{-k}\).

We next express the information difference in terms of exact projection
densities.  Let \(Q\) denote the independent
law preserving \((S,W)\).  Their \((D,S)\) and \((G,S)\)
marginals agree after identifying the matrices.  Hence
\(h_Q(GS\mid G)=h_P(DS\mid D)\).
Proposition~\ref{prop:r-information-comparison} gives
\begin{equation}\label{rlevels:two-information-identity}
 \mathcal J(W)-I_P(S;W\mid D,DS)
 =h_Q(GS\mid G,W)-h_P(DS\mid D,W)-I_P(S;A\mid W).
\end{equation}
Here independence of \(B\) from \((S,A,W)\) reduces
\(I_P(S;D\mid W)\) to \(I_P(S;A\mid W)\).
Lemma~\ref{lem:r-projection-density} makes the baseline
and conditional entropies finite; it also gives the actual conditional
image density.  Thus \eqref{rlevels:two-information-identity} is an
identity of finite quantities, with the displayed negative sign.

For a matrix \(D\), let \(p_{w,D}\) be a density of the image of
\(\nu_w\) under \(s\mapsto Ds\).  This denominator uses
\(\nu_w=P_{S\mid W=w}\) before conditioning on the actual rows.
Contraction of
\[
 \mathbb E D(P_{S\mid A,W}\Vert\nu_W)=I_P(S;A\mid W)
\]
under \(s\mapsto Ds\), with \(B\) still independent, gives
\begin{equation}\label{rlevels:two-cross-entropy}
 -h_P(DS\mid D,W)-I_P(S;A\mid W)
 \le \mathbb E_P\log p_{W,D}(DS).
\end{equation}
The numerator after contraction is the true conditional image density
given \((D,W)\).  Its log ratio to \(p_{W,D}\) is integrable:
the relative entropy is finite, and the negative part of a log
likelihood ratio has expectation at most \(1/e\).  Subtracting this
ratio from the integrable actual log density also proves integrability
of the right side.

Here we choose a density version that retains a bound at every fixed
target \(s\).  For each positive integer \(a\), let
\(\varphi_\delta^a\) be the centered Gaussian density of covariance
\(\delta^2I_a\), and fix a sequence \(\delta_t\downarrow0\).
For any probability measure \(\nu\le K\sigma\), define
\[
 F_\delta^\nu(g,s)
   =\int\varphi_\delta^k(g(u-s))\,\nu(du),\qquad
 p^\flat_\nu(g,y)
   =\liminf_{t\to\infty}
       \int\varphi_{\delta_t}^k(gu-y)\,\nu(du).
\]
These functions are measurable.  For each full-rank \(g\), the
Gaussian approximate identity converges to the image density at its
Lebesgue points, so \(p^\flat_\nu(g,\cdot)\) is a density version.
It may be defined by a liminf on the exceptional set; this choice is
useful at a fixed target.  It is valid in
\eqref{rlevels:two-cross-entropy}, since the actual conditional image
law is absolutely continuous.

There is a quantitative first-density passage under the independent
law, before taking logarithms.  For independent \(G\sim\gamma_k\)
and \(S\sim\nu\), Gaussian integration and Tonelli give
\begin{align}
 \mathbb E_{G,S}F_\delta^\nu(G,S)
 &= (2\pi)^{-k/2}\iint
       (\|u-s\|^2+\delta^2)^{-k/2}\,\nu(du)\nu(ds) \notag\\
 &\le (2\pi)^{-k/2}\iint
       \|u-s\|^{-k}\,\nu(du)\nu(ds)<\infty.
 \label{rlevels:two-first-density}
\end{align}
The last integral is finite uniformly in \(s\): the diagonal has
zero mass, and the centered cap bound on dyadic shells gives
\[
 \sup_{s\in S^{d-1}}\int\|u-s\|^{-k}\,\nu(du)
 \le 1+K2^k\sum_{i\ge0}2^{-i(N-k)}<\infty.
\]
At \(GS\) the approximate identity converges almost surely under
the independent law, by absolute continuity of the image, Lebesgue
differentiation, and Fubini.  Fatou in
\eqref{rlevels:two-first-density} therefore gives
\[
 \mathbb E_{G,S}p^\flat_\nu(G,GS)
 \le (2\pi)^{-k/2}\iint\|u-s\|^{-k}\,\nu(du)\nu(ds)<\infty.
\]
This proves the positive logarithmic part directly.  For the negative
part compare with the standard Gaussian density \(\phi\) on
\(\mathbb R^k\): \(t(\log t)_-\le1/e\) bounds the negative part
of \(\log(p^\flat_\nu/\phi)\), and
\(\mathbb E\|GS\|^2=k\) bounds the cross entropy with \(\phi\).
This recovers independent log integrability for this density version.
Actual log integrability is supplied by the contraction preceding
\eqref{rlevels:two-cross-entropy}.  In what follows take
\(p_{w,D}=p^\flat_{\nu_w}(D,\cdot)\).  The mollifiers here only
specify and estimate densities at exact labels.

It remains to bound those densities under independent rows before
changing the law of \(A\).  We claim that for every fixed
\(w\) of positive probability and every \(s\in S^{d-1}\),
\begin{equation}\label{rlevels:two-moment}
 \left[\mathbb E_A
 \bigl(\mathbb E_B p^\flat_{\nu_w}((A;B),(A;B)s)\bigr)^q
 \right]^{1/q}
 \le e^{Cd}U_w(1+\log\Gamma_w),
\end{equation}
where \(A\sim\gamma_n\) and \(B\sim\gamma_l\) are
independent in this displayed expectation.  We prove the
uniform-in-\(\delta\) version first, with
\[
 F_\delta(A,B,s)
   =\int\varphi_\delta^k((A;B)(u-s))\,\nu_w(du).
\]
For \(h=u-s\),
\[
 \mathbb E_B\varphi_\delta^l(Bh)
 =(2\pi)^{-l/2}(\|h\|^2+\delta^2)^{-l/2}.
\]
Expanding the \(q\)-th power introduces \(u_1,\ldots,u_q\) and
\(h_i=u_i-s\).  For their Gram matrix
\(K=(\langle h_i,h_j\rangle)_{i,j}\), Gaussian integration in the
\(n\) shared rows gives exactly
\[
 \mathbb E_A\prod_{i=1}^q\varphi_\delta^n(Ah_i)
 =(2\pi)^{-nq/2}\det(K+\delta^2I_q)^{-n/2}.
\]
Using \(\det(K+\delta^2I_q)\ge\det K\) and Gram--Schmidt,
we conclude
\begin{equation}\label{rlevels:two-gram}
 \mathbb E_A\bigl(\mathbb E_B F_\delta(A,B,s)\bigr)^q
 \le (2\pi)^{-kq/2}
 \int\prod_{i=1}^q
   \|u_i-s\|^{-l}\operatorname{dist}(u_i-s,L_{i-1})^{-n}
   \,\nu_w^{\otimes q}(d\mathbf u),
\end{equation}
where \(L_{i-1}\) is the span of \(u_1-s,\ldots,u_{i-1}-s\),
starting with \(\{0\}\).  An infinite upper bound is allowed at
singular Gram matrices.  All integrations so far are nonnegative.

We bound the corresponding one-point integral uniformly over \(s\)
and every linear subspace \(L\) of dimension at most \(q\).
On a radial shell \(R/2<\|u-s\|\le R\), \(R=2^{1-i}\),
the part at distance at most \(2^{-a}R\) from \(L\) is covered
by at most \((3\cdot2^a)^q\) balls of radius \(2^{1-a}R\).
This follows by covering the radius-\(R\) projection ball in \(L\)
with a \(2^{-a}R\)-net.  By \eqref{rlevels:two-two-powers} its
\(\nu_w\)-mass is at most
\[
 C^d\min\{
 U_w2^{-a(k-q)}R^k,\,
 \Gamma_w2^{-a(N-q)}R^N\}.
\]
In particular the set at distance zero has zero mass.  On the band
\(2^{-a-1}R<\operatorname{dist}(u-s,L)\le2^{-a}R\), the
integrand \(\|u-s\|^{-l}\operatorname{dist}(u-s,L)^{-n}\)
is at most \(2^k2^{an}R^{-k}\).  Sum the transverse bands first.
The positive gaps \(k-q-n\) and \(N-q-n\) make both sums geometric,
so the integral on the radial shell is at most
\[
 C^d\min\{U_w,\Gamma_wR^{N-k}\}.
\]
For the radial sum put \(a_0=N-k\ge d/3\).  The sum of
\(\min\{1,(\Gamma_w/U_w)2^{a_0(1-i)}\}\) has at most
\(2+(\log_2(\Gamma_w/U_w))_+/a_0\) saturated terms followed by
a geometric tail.  Since \(U_w\ge2^{-k}\) and
\(\Gamma_w\ge1\), this proves
\begin{equation}\label{rlevels:two-inverse}
 \int\|u-s\|^{-l}\operatorname{dist}(u-s,L)^{-n}\,\nu_w(du)
 \le e^{Cd}U_w(1+\log\Gamma_w).
\end{equation}
Integrate \eqref{rlevels:two-gram} from its last point to its first,
applying \eqref{rlevels:two-inverse} each time.  This proves the
claimed bound uniformly in \(\delta\).

The liminf version now gives the exact fixed-target conclusion:
for every fixed \(s\),
\[
 \mathbb E_Bp^\flat_{\nu_w}((A;B),(A;B)s)
 \le\liminf_{t\to\infty}\mathbb E_BF_{\delta_t}(A,B,s)
\]
by Fatou in \(B\).  Raise to the \(q\)-th power and apply Fatou
in \(A\).  The uniform moment bound proves
\eqref{rlevels:two-moment}, including at targets where a density
version specified only almost everywhere would not have sufficed.
This is the distinct fixed-target passage used in this comparison.

We finally pay for dependence of the actual rows.  Their conditional
relative entropies satisfy
\begin{equation}\label{rlevels:two-row-kl}
 \mathbb E D(P_{A\mid S,W}\Vert\gamma_n)
 =I_P(A;W\mid S)=I_P(A;V\mid S)\le H(V).
\end{equation}
Here \(A,S\) are marginally independent, and \(W\) adds only
a function of \((S,V)\).  Given \((S,W,A)\), \(B\) is still
independent.  Jensen in \(B\), followed by
\eqref{rlevels:tilt} for the true \(P_{A\mid S,W}\) relative to
\(\gamma_n\), uses \eqref{rlevels:two-moment} to give
\begin{equation}\label{rlevels:two-actual-upper}
 \mathbb E_P\log p^\flat_{\nu_W}(D,DS)
 \le \frac{H(V)}q+Cd+
 \mathbb E\bigl[\log\alpha_W+b\log2+
                    \log(1+\log\Gamma_W)\bigr].
\end{equation}
Finite average divergence ensures the required absolute continuity.
Truncation permits the conditional inequality wherever a log
expectation is negative infinite; the integrated cross entropy is
finite by \eqref{rlevels:two-cross-entropy}.

The \(b\log2\) term cancels
\eqref{rlevels:two-independent-entropy} when substituted with
\eqref{rlevels:two-actual-upper} into
\eqref{rlevels:two-information-identity}--\eqref{rlevels:two-cross-entropy}.
Moreover
\[
 \mathbb E\log\alpha_W=H(W\mid V),\qquad
 \mathbb E\log(1+\log\Gamma_W)\le\log(1+H(W)).
\]
Removing the refinement costs another \(H(W\mid V)\), since
\[
 \mathcal J(V)\le\mathcal J(W),\qquad
 I_P(S;W\mid D,DS)\le I_P(S;V\mid D,DS)+H(W\mid V).
\]
Thus the full quantitative budget is
\begin{equation}\label{rlevels:two-budget}
 \mathcal J(V)-I_P(S;V\mid D,DS)
 \le Cd+\frac{H(V)}q+2H(W\mid V)+\log(1+H(W)).
\end{equation}
Only now use \(H(V)\le d^2\), \(q=\lfloor d/8\rfloor\), and
\eqref{rlevels:two-label-entropy}.  They bound the right side by
\(Cd\), proving \eqref{rlevels:two-comparison-bound}.
\end{proof}


The applications of both comparisons, including their distinct residual-sphere constants, are proved in Subsection~\ref{rlevels:endpoints}.
